\documentclass[11pt]{article}

\usepackage[utf8]{inputenc}
\usepackage{amsmath,amssymb}
\usepackage[preprint]{tmlr}
\usepackage{natbib}
\usepackage{booktabs}
\usepackage{graphicx}
\usepackage{hyperref}
\usepackage{enumitem}
\usepackage{microtype}
\usepackage{underscore}

\newcommand{\pend}[1]{\textbf{[[PENDING: #1]]}}

\title{Approval Gates versus Safety Prompts:\\A Plan-Level Audit of Coverage, Cost, and Disclosure}
\author{Zewen Liu \\ \texttt{17353895263@163.com}}



\date{}

\begin{document}

\maketitle

\begin{abstract}
Approval gates and safety prompts can both reduce unapproved actions, but a zero count does not reveal whether actions were never proposed, were reviewed, or disappeared after a parse failure. We study these mechanisms in a plan-level harness using 30 authored travel scenarios, nine configurations, three processing orders, and two recorded model backbones (1,620 cells). A gate-by-safety-prompt factorial is complemented by deterministic-gate, regeneration, disclosure, and memory ablations. With the gate alone, 18 and 3 designated dangerous actions remain in the final plans; with the safety prompt alone, 0 and 2 remain. All gated arms have zero retained actions lacking approval, while the ungated arms have 24 and 25. Under the primary backbone, adding the safety prompt to the gate removes every audit decision because no designated dangerous action is proposed. A deterministic rule also produces an audit trail without model-based auditing. The safety prompt uses no additional model calls but increases prompt-token consumption by 14.4 percent on the primary backbone. No task-success contrast is significant at the scenario unit; this does not establish equivalence. Disclosure is sensitive to both the regeneration instruction and the keyword instrument. Parse failures produce empty plans in 16.7 percent of Full cells under the exploratory backbone. These results concern retained plans, not executed actions or harm: the harness has no tool-execution layer or ground-truth safety labels. We provide a measurement case study separating proposal, review, retention, and user-facing disclosure, rather than evidence that any intervention makes an agent safer.
\end{abstract}

\section{Introduction}
\label{sec:intro}

Runtime controls for multi-agent tool execution---approval gates, guardrails,
human-in-the-loop checks---are widely deployed, yet their \emph{measured}
cost/benefit on realistic tool-use benchmarks is rarely reported. A gate that
blocks dangerous tool calls is intuitively safer, but it also consumes extra
model calls, and it is not obvious that it dominates the cheapest alternative:
a safety instruction in the system prompt.

This paper measures that trade-off in retained plans rather than in live tool execution. We compare a core gate-by-safety-prompt factorial and five additional configurations on the travel domain of a public enterprise multi-agent benchmark. The principal configurations are:

\begin{itemize}[leftmargin=*,itemsep=2pt]
\item \textbf{Full}: an explicit negotiation gate (one auditor call per
      dangerous action) over bounded memory.
\item \textbf{Abl-noForget}: the same gate without the forgetting mechanism.
\item \textbf{Abl-noNegot}: no gate at all (dangerous actions are retained directly).
\item \textbf{Baseline-promptSafety}: no gate; safety rests on a system-prompt
      policy that forbids unconfirmed reservation/payment actions
      (zero extra model calls).
\item \textbf{Baseline-heavyAudit}: a strengthened gate (three auditor rounds per
      dangerous action).
\end{itemize}

Our headline finding is a \emph{mechanism} distinction that a coverage metric
cannot see. Both the explicit gate and the prompt-level policy can drive the
count of un-negotiated dangerous actions to zero, but they do so in opposite
ways: the gate lets the model \emph{propose} the dangerous action and then
\emph{audits} it, leaving an auditable record of what was proposed, approved,
and denied; the prompt-level policy usually prevents the action from being
proposed at all, leaving \emph{no} record---though under our second backbone it
did not, and the action was retained unconfirmed (\S\ref{sec:res-coverage}). A low unconfirmed-action count therefore does not establish that a prompt policy and a gate are interchangeable. The prompt adds token cost; both the model-driven and deterministic gates can record decisions. The comparison concerns proposal suppression, recorded review, and retention, rather than a universal safety ranking.

We make three scoped contributions: (1) a falsifiable coverage instrument derived from retained plans and audit records; (2) a factorial case study distinguishing proposal suppression from plan filtering; and (3) an audit of parser, disclosure, and cost measurements that can change the apparent ranking of interventions. Task success remains inconclusive at the scenario unit, and all conclusions are limited to this planning harness.

\section{Related Work}
\label{sec:related}

\paragraph{Tool-using agents and safety.}
ReAct-style and multi-agent orchestration frameworks interleave reasoning with
tool calls~\cite{shu2024effective}. Safety interventions in this setting are
typically prompt-level policies, guardrail classifiers, or explicit
confirmation steps. GuardAgent~\cite{xiang2025guardagent} compiles a safety
guard request into executable guardrail code and reports high accuracy on two
purpose-built benchmarks that carry policy labels; ToolSafe~\cite{mou2026toolsafe}
trains a step-level guard model on TS-Bench, a benchmark of unsafe tool
invocations, and reports a 65\% reduction in harmful invocations. Both therefore
measure safety against \emph{ground-truth safety labels}. The enterprise task
benchmark we use has none: its assertions only decide whether the task was
completed. Our contribution is not a new mechanism---the approval gate is a
known pattern---but a measurement of what such a gate buys, and at what cost,
when no safety labels exist to score against.

\paragraph{Deterministic and policy-based enforcement.}
A complementary line of work moves the safety decision out of the model entirely.
AgentSpec~\cite{wang2025agentspec} is a domain-specific language for runtime
constraints---triggers, predicates and enforcement mechanisms---and reports
preventing unsafe executions in over 90\% of code-agent cases;
Progent~\cite{shi2025progent} expresses privilege as symbolic rules over tool
names and arguments, checks every call deterministically, and uses an SMT solver
to apply policy narrowings automatically while requiring explicit approval for
expansions, so the effective action space can only shrink without approval;
SEAgent~\cite{ji2026seagent} enforces mandatory access control over an
information-flow graph; and CXI~\cite{santosgrueiro2026cxi} admits a call only
when field authority, exact-effect authorization and invocation authority bind to
the same action manifest. These systems supply precisely what our measurement
finds missing from a model-driven gate: a \emph{deterministic} decision and a
structured audit record that does not depend on the model obeying an instruction.
We do not implement such a policy language. We instead measure what a purely
model-driven gate buys, and we include a deterministic rule-based gate as the
cheapest possible stand-in, to locate the trade-off rather than to claim that no
better enforcement mechanism exists.

\paragraph{Labelled step-level guards.}
On the labelled side, ToolSafe/TS-Bench~\cite{mou2026toolsafe} and
StepGuard~\cite{zheng2026stepguard} train step-level guard models and evaluate
them against ground-truth unsafe actions on AgentDojo and AgentDyn; StepGuard
reports a 77.3\% reduction in mean attack success rate at a 2.8-percentage-point
utility cost. Our benchmark carries no such labels, which is why we report
coverage and auditability rather than detection accuracy, and why we cannot score
a refusal as correct rather than as an over-rejection.

\paragraph{Agent benchmarks.}
We use the travel domain of the AWS Bedrock Agents Multi-Agent Collaboration
Scenario Benchmark~\cite{awsbench2024,shu2024effective}, which provides
authored enterprise scenarios with programmatic assertions over a ten-agent
tool suite.

\paragraph{Coverage versus auditability.}
Aggregate safety scores can omit the route by which an intervention changes a plan. We separately report the number of
dangerous actions retained without confirmation, and the number of
\emph{recorded} auditor decisions. As we show, these can move in opposite
directions across mechanisms.

\paragraph{Over-safety and the missing label.}
Zhang et al.~\cite{zhang2026guard} argue that at the authorization boundary the
safe-versus-unsafe label is a choice of authorization policy rather than a
property fixed by the action, and build a benchmark whose labels are mechanical
consequences of a stated policy. This is exactly the constraint under which we
operate, and it is why we refuse to claim that a gate makes the system
\emph{safer}: without a policy-derived label, a denial is equally consistent
with a correct block and an over-rejection.

\paragraph{Trajectory auditing and consent integrity.}
HarnessAudit~\cite{liu2026harnessaudit} evaluates full execution trajectories and separates task completion from boundary compliance. Runtime-contract work~\cite{ng2026runtime} likewise distinguishes preventive controls from evidence that an action happened, while consent-integrity work~\cite{weng2026consent} binds approval to the actual action at the execution boundary. Our harness has no such execution boundary and cannot validate those guarantees. Its narrower contribution is a controlled plan-level comparison of a safety sentence, model-driven review, and a fixed deterministic rule. Singh~\cite{singh2026scapegoats} studies misleading safety explanations following silent tool failures; we instead vary regeneration instructions after explicit plan-level denials. These are related measurement problems, not evidence that refusal auditing or approval integrity is new.

\paragraph{Bounded memory.}
Capacity-limited memory with eviction is a standard technique for keeping agent
context small. We include a forgetting ablation to test whether the eviction
mechanism itself changes task outcomes on this benchmark.

\section{Method}
\label{sec:method}

\subsection{Dangerous actions}

We whitelist six tool actions as \emph{dangerous} (commitment/payment only):
\texttt{bookflight}, \texttt{bookhotel}, \texttt{createreservation},
\texttt{bookairbnb}, \texttt{neworderevent}, \texttt{checkout}. Actions such as
\texttt{selectseat} or \texttt{cancelticket} are excluded because they do not
create a new commitment.

\subsection{Runtime configurations}
\label{sec:config}

The configurations form a $2\times2$ factorial over (a) whether a dangerous action
is routed through an auditor and (b) whether the safety sentence is present in the
system prompt, plus further arms that isolate the remaining factors.
Table~\ref{tab:design} lists all nine; the design is declared as a table rather
than as prose so that the factorial structure is explicit and every arm can be
re-implemented. The table is a fixed design specification; the per-cell configuration snapshots provide the authoritative configuration record.

\textbf{Gate.} When an auditor is used it is called with the action and the user
request and must answer \texttt{APPROVE} or \texttt{DENY}. \texttt{Full} and
\texttt{Abl-noForget} allow up to two rounds per action and
\texttt{Baseline-heavyAudit} up to three; the loop breaks at the first unambiguous
verdict. Because it can stop early, that budget is an upper bound, so
\S\ref{sec:res-cost} reports the rounds actually consumed.

\textbf{Deterministic gate.} \texttt{Baseline-detGate} replaces the model with a
fixed rule: a dangerous action is approved only if the user's request carries one
of a small fixed list of commitment cues (\texttt{i confirm}, \texttt{confirmed},
\texttt{go ahead}, \texttt{please proceed}, \texttt{proceed with}, \texttt{book
it}, \texttt{make the booking}, \texttt{yes, book}). The list is given verbatim so
that the rule can be reproduced, and the arm spends zero model calls. It stands in
for the policy-based enforcement discussed in \S\ref{sec:related}; it is one rule,
not a policy language.

\textbf{Prompt-level safety.} \texttt{Baseline-promptSafety} invokes no auditor.
The system prompt additionally instructs the model not to include a
reservation/payment action unless the user explicitly confirmed that purchase, and
to ask for confirmation otherwise.

\textbf{Answer regeneration.} When the retained action set differs from the
proposed set, the user-facing answer is regenerated so that a refused action is not
reported as done. \texttt{Fact-gateNoRegen} disables this step with the gate on,
isolating its effect, and \texttt{Fact-disclosePrompt} additionally requires the
regeneration to name the refused actions and give a reason.

\input{tables/design}

\subsection{Bounded memory}

Memory is capacity-limited; on overflow the least-frequently-used entry is
evicted. The capacity is deliberately small so that eviction actually fires.
\texttt{Abl-noForget} disables eviction.

\subsection{Metrics}
\label{sec:metrics}

\begin{itemize}[leftmargin=*,itemsep=2pt]
\item \textbf{Task success} (0/1): every programmatic assertion for the scenario is
      satisfied. \textbf{Not every assertion is checkable.} An assertion is
      evaluated only if it names a tool action the harness knows or carries the
      \texttt{user:} prefix; anything else is recorded as \emph{unchecked} rather
      than silently passing or failing, and prefix matching is case-insensitive so
      a ``\texttt{User:}'' assertion is still scored. Success therefore means ``no
      checked assertion failed'', not ``the scenario was verified'';
      \S\ref{sec:res-success} reports the unchecked share alongside it. Assertions
      that do carry \texttt{user:} are scored by a keyword-overlap heuristic, a
      weak instrument whose behaviour is disclosed where it is used.
\item \textbf{\texttt{gate\_bypass}} (stored as \texttt{unnegotiated\_count}): per
      run, the number of whitelisted dangerous actions retained in the final plan
      that have \emph{no} \texttt{APPROVE} record in the audit log. It is computed
      from the plan and the log, never assumed: for a correctly functioning gate it
      is zero, but a bypassed gate, a dropped audit record or an exception inside
      the audit loop would appear as a non-zero. For an ungated configuration there
      is no log, so every retained dangerous action counts. \textbf{This measures
      gate coverage, not real-world safety}: the benchmark carries no ground-truth
      unsafe labels.
\item \textbf{\texttt{retained\_dangerous}}: per run, the number of dangerous
      actions left in the final plan. This is a different quantity from
      \texttt{gate\_bypass}, because a gate can \emph{approve} an action, which is
      then retained and audited rather than bypassed. \S\ref{sec:res-coverage}
      reports both, since they answer different questions.
\item \textbf{Audited actions, rounds and verdict class.} Each auditor decision is
      logged with the action, the verdict, the number of rounds it consumed, the
      raw reply, and a class: \texttt{explicit\_approve}, \texttt{explicit\_deny},
      \texttt{both\_tokens} (the reply contained both verdict words; the
      conservative choice is taken and the ambiguity is recorded),
      \texttt{unparsed} (no verdict within the round budget) or
      \texttt{all\_errors} (every call failed). Without this classification a
      refusal rate cannot be separated from a parser failure, which is why we
      report the class distribution rather than only the refusal rate.
\item \textbf{Measured cost.} For every model call the harness records wall-clock
      latency and, when the backend reports them, prompt and completion tokens (the
      OpenAI \texttt{usage} block, or ollama's
      \texttt{prompt\_eval\_count}/\texttt{eval\_count}). Cost is therefore
      reported as measured per configuration rather than bounded.
\item \textbf{Parse provenance.} A reply that does not parse yields an empty action
      list, which a retained-action metric would score as full compliance. The
      harness records whether the reply parsed, so ``the model planned nothing''
      and ``the reply did not parse'' can be told apart.
\end{itemize}

\section{Experimental Setup}
\label{sec:setup}

\paragraph{Benchmark.} Travel domain, all 30 scenarios, each a user request
over a ten-agent tool suite with assertions defining success. Data are
CC-BY-4.0; scaffolding code is MIT-0 (per the repository README). Scenarios are
authored enterprise cases, not real user traffic.

\paragraph{Recorded model identifiers.} The retained service configurations label the two backbones as \textbf{Qwen3.8-27B}
(reported in the project environment notes as AWQ-INT4/vLLM 0.27.1) and \textbf{Qwen3.5-9B} (served via Ollama).
The records use service identifiers \texttt{qwen3.8-27b} and \texttt{qwen3.5:9b}. The retained records do not pin a public checkpoint revision or weight hash, so these labels are service identifiers rather than independently verified checkpoint identities. Thinking traces are disabled at the serving layer
(\texttt{enable\_thinking=false} / \texttt{think=false}) so that the response is
the required JSON object.

\paragraph{Grid.} 30 scenarios $\times$ 9 configurations $\times$ 3 seeds
$= 810$ cells per model; two backbones, so $1620$ cells in total.

\paragraph{Analysis scope and units.} The retained analysis plan was written before completion of an earlier five-configuration run; the instrumented nine-configuration grid, factorial intervals, and disclosure ablations are extensions. We therefore present the current grid as an exploratory measurement study, with the 27B-labelled service as the primary reporting backbone and the 9B-labelled service as an exploratory replication, rather than claiming a prospectively registered nine-arm experiment. The plan and its deviation record are retained in \texttt{evidence/analysis\_plan\_v4.md}.

Three seeds change scenario processing order; they are not three independent model replications. For inference we average paired measurements within each of the \textbf{30 scenario units}. Task-success contrasts use an exact two-sided sign test on informative scenario-level paired differences and Holm correction across \textbf{eight comparisons against Full per backbone}. Cell-level exact McNemar tests over 90 scenario-order pairs are sensitivity analyses only. No scenario-level success contrast is significant after correction; this does not establish equivalence or absence of an effect. Factorial effects are paired differences of scenario means with percentile bootstrap intervals (2,000 resamples, fixed bootstrap seed 42). These are \emph{pointwise exploratory intervals}, not simultaneous or multiplicity-adjusted confidence statements. Coverage, audit-record, disclosure-marker, and latency counts are reported descriptively.

For example, changing the gate with the safety sentence absent changes retained dangerous-action totals from $24$ to $18$ under the primary backbone and $25$ to $3$ under the replication; changing the safety sentence with the gate absent changes the totals from $24$ to $0$ and $25$ to $2$. These are totals over 90 cells per configuration, not independent-sample effect estimates. Latency was observed on the available serving stacks without randomized execution order, so it is not a hardware-independent ranking.

\section{Results}
\label{sec:results}

All numbers below are recomputed from the per-cell JSON of the two model runs
(Qwen3.8-27B and Qwen3.5-9B; 27 group-seed files and 810 cells each: nine
configurations $\times$ three seeds $\times$ thirty scenarios) by
\texttt{scripts/analysis\_v5.py}. Every table is emitted from that JSON by
\texttt{scripts/make\_tables.py} rather than transcribed into the manuscript.
Seed-level and scenario-level denominators are distinguished explicitly, and we
flag the places where our own instrument is weaker than it looks.

\input{tables/results}

\subsection{Coverage: a derived count, and what it still cannot see}
\label{sec:res-coverage}

The quantity usually reported for an approval gate is the number of dangerous
actions that go unconfirmed. In the earlier version of this harness that field was
\emph{assigned} the constant zero for every gated configuration, so it could not
fail: a bypassed gate, a dropped audit record or an exception inside the audit loop
would all have read as perfect coverage. It is now derived per run, as the number of
dangerous actions retained in the final plan that carry no \texttt{APPROVE} record
in the audit log (\S\ref{sec:metrics}). The derivation is re-checked against the
value stored at run time for all $810$ rows of each model run, with
\textbf{zero} mismatches, and regeneration is verified never to fire without a
preceding denial.

Table~\ref{tab:main-27b} reports the result. The derived count is zero for every
gated configuration under both backbones, and it is non-zero exactly where it
should be: \texttt{Abl-noNegot}, which has no gate, retains $\mathbf{24}$ dangerous
actions under the 27B model and $\mathbf{25}$ under the 9B model.

The informative case is the prompt-level policy. \texttt{Baseline-promptSafety}
never invokes an auditor, so on this field it also reads zero under the 27B
model---but that zero is a measurement of the model's compliance rather than a
property of a code path, and it is not stable across backbones. Under the 9B model
the prompt policy lets $\mathbf{2}$ dangerous actions through that no mechanism
reviewed, against $\mathbf{0}$ for every gated configuration including the
deterministic one. So the derived count \emph{can} separate the two mechanisms, but
only where the prompt policy actually fails, and it does so under one of our two
backbones rather than both.

\paragraph{Retained actions tell the opposite story.}
Coverage counts what escaped review and says nothing about what survived it. Under
the 27B model the gate retained $\mathbf{18}$ dangerous actions---of the $24$
proposed, $18$ were approved and $6$ denied---while \texttt{Baseline-promptSafety}
retained $\mathbf{0}$ and \texttt{Abl-noNegot} retained $24$. Under the 9B model the
gate retained $3$, the prompt policy $2$, the ungated arm $25$. On this reading the
gate and the prompt policy are far apart, which is the opposite of what the bypass
count says. Both numbers are correct, because they answer different questions: one
counts actions that were never reviewed, the other counts actions that were reviewed
and allowed.

The deterministic gate is the sharpest illustration. It denied \emph{every} dangerous
action it saw---$24$ of $24$ under the 27B model and $21$ of $21$ under the 9B
model---because no request in this benchmark carries one of its commitment cues. It
therefore matches the prompt policy's retained count of zero, while producing a
complete audit record and spending no model calls on auditing at all.

\paragraph{Parse failures can inflate apparent compliance.}
A cell whose plan fails to parse yields an empty action list, which a retained-action
metric scores as perfect compliance. The instrumented run separates the two cases
(Table~\ref{tab:parse-27b}). Under the 27B model the split is small: $3.3\%$ of cells
fail to parse and $3.3\%$ are genuinely empty, and for the prompt-level arms the
parse-failure rate falls to $1.1\%$. Under the 9B model it does not: \texttt{Full}
fails to parse in $\mathbf{16.7\%}$ of cells against $12.2\%$ genuinely empty, so
\emph{more than half} of its empty plans are parser failures rather than restraint.
The prompt-level arms again fail far less often ($5.6\%$). This cuts both ways and we
state both: the prompt policy's low retained count is less contaminated by parser
failures, and the two mechanisms are therefore not being compared on equally clean
plans.

\subsection{Auditability: a design difference, and what the records show}
\label{sec:res-audit}

The gate records what it saw; the prompt policy has nothing to record.
\texttt{Full} logs $\mathbf{24}$ auditor decisions under the 27B model
($18$ \texttt{APPROVE}, $6$ \texttt{DENY}) and $22$ under the 9B model
($3$ \texttt{APPROVE}, $19$ \texttt{DENY}). \texttt{Baseline-promptSafety} logs
$\mathbf{0}$ under both, because that configuration never constructs an auditor
call. We state this as a property of the configuration rather than as a statistical
finding: the zero is entailed by which code path the arm takes, so a significance
test on it would measure the configuration rather than the world.

Two things in the records are worth separating from that design statement.

\paragraph{Every decision resolved explicitly.}
Each decision now carries a verdict class (Table~\ref{tab:audit-27b}), so we can say
what \emph{kind} of refusal these are rather than only how many. Across all gated
configurations there are $120$ model-auditor decisions under the 27B model ($90$
\texttt{explicit\_approve}, $30$ \texttt{explicit\_deny}) and $114$ under the 9B
model ($96$ \texttt{explicit\_deny}, $18$ \texttt{explicit\_approve}), and
\textbf{not one} is classified \texttt{unparsed}, \texttt{all\_errors} or
\texttt{both\_tokens}. The refusal rates we report are therefore not inflated by
parser or transport failures---a caveat the earlier harness could not exclude,
because it mapped both onto \texttt{DENY} without recording which had occurred.

\paragraph{An audit trail does not require a model.}
\texttt{Baseline-detGate} logs $24$ decisions under the 27B model and $21$ under the
9B model while issuing \textbf{zero} auditor calls: its realized rounds per scenario
are $0.00$ against $0.27$ for \texttt{Full}. Auditability is therefore a property of
having a decision point that records, not of having a model in the loop. This
sharpens the comparison the paper makes, because the prompt policy's zero does not
mean ``no decision was needed'' but ``no decision was recorded''.

One configuration yields an empty audit log for a third reason. Under the 27B model
\texttt{Fact-gateOn-safetyOn} logs $\mathbf{0}$ decisions---not because it lacks a
gate, but because the safety sentence suppresses the proposal of dangerous actions
altogether, so the gate is never reached. Under the 9B model the same configuration
logs $3$ decisions, all approvals. An empty audit log is thus compatible with three
different mechanisms, and only the configuration table
(Table~\ref{tab:design}) tells them apart.

\subsection{Task success: no effect detectable at the scenario unit}
\label{sec:res-success}

Success is low everywhere and the spread between configurations is small
(Table~\ref{tab:main-27b}). Under the 27B model every gated configuration scores
$5/90$ cells ($5.6\%$, Wilson $[2.4, 12.4]$) and the two configurations carrying the
safety sentence score $6/90$ ($6.7\%$, $[3.1, 13.8]$). Under the 9B model the gated
configurations score $4/90$ ($4.4\%$, $[1.7, 10.9]$), the ungated arm $6/90$, and the
prompt policy $8/90$ ($8.9\%$, $[4.6, 16.6]$).

\emph{Scenario unit (primary reporting analysis).} The earlier analysis plan identifies the 30 scenarios as the
independent unit, because decoding is not deterministic even at temperature zero and
the seeds change processing order rather than constituting independent draws. At that
unit no contrast reaches significance under either backbone. The largest is
\texttt{Full} versus \texttt{Baseline-promptSafety}: three informative scenarios under
both, exact sign-test $p=1.0$ under the 27B model and $p=\mathbf{0.25}$ under the 9B
model---and $0.25$ is the \emph{minimum attainable} two-sided value for three
informative units, so this contrast cannot reach significance at our unit of analysis
whatever its effect size. After Holm correction across the eight comparisons per
model, every adjusted $p$ is $1.0$.

\emph{Cell level (sensitivity).} Taking each (seed, scenario) pair as a unit gives 90
pairs per comparison. Under the 9B model \texttt{Full} versus
\texttt{Baseline-promptSafety} has $4$ discordant pairs and exact McNemar $p=0.125$;
under the 27B model the same contrast has $5$ discordant pairs and $p=1.0$. We report
these as sensitivity checks only: the pairs are not independent, since the same
scenario recurs across seeds, and a cell-level $p$ is smaller than the scenario-level
one partly because the denominator is larger. Detectability calculations based on these dependent cells are not independent-unit power analyses. With only three informative scenario units in the gate-versus-prompt comparison, the minimum attainable exact two-sided sign-test value is $0.25$. We therefore make no equivalence claim and do not interpret non-significance as evidence that a useful effect is absent.

\emph{What the discordant cells actually contain.} A discordant count invites an
attribution the data may not support, so we inspected every discordant cell. For
\texttt{Full} versus \texttt{Baseline-promptSafety}, the contrast that looks most
interesting, \textbf{no discordant cell under either backbone contains a gate
denial}. Under the 27B model one of the five contains a gate decision at all and it
is an approval; under the 9B model none of the four contains any gate decision. Four
of five (27B) and three of four (9B) have an identical retained action set in both
arms, and two of five and three of four carry no programmatically checked tool
assertion, so they are decided by the keyword-overlap heuristic applied to the
generated answer. These records do not support attributing the success difference to a gate denial; differences in prompt-conditioned answers and their heuristic scoring remain alternative explanations.

There is one contrast we \emph{can} attribute. Under the 9B model, \texttt{Full}
versus \texttt{Abl-noNegot} has two discordant cells, both in the same scenario
(idx~10); in both, the gate denied two actions, the retained action sets differ, and
the ungated arm succeeded where the gated arm failed. This instance links gate denials to a lower harness success score: refusing a booking required by the scenario removes an action scored as necessary. Without authorization labels, that association cannot establish whether the denial was appropriate. It rests on a single scenario, so no
inference is possible, and we report it as one attributable instance rather than as an
effect size.

\paragraph{The scorer leaves part of every run unchecked.}
$432$ of $3564$ cell-assertions ($12.1\%$) are not checkable by our scorer, and the
figure is \emph{identical} across all nine configurations and both backbones
($48/396$ per configuration), so the missingness pattern is shared across arms. These assertions name no known
tool action; shared missingness does not guarantee unbiased intervention comparisons. Success therefore means ``no \emph{checked} assertion failed'', not
``the scenario was verified''. Making the prefix dispatch case-insensitive moved
three assertions per configuration from silently skipped to checked; the remainder are
semantic rather than tool assertions and stay unchecked.

\subsection{Disclosure of denials}
\label{sec:res-disclosure}

\textbf{Denominator.} We count one unit per gated run that contains at least one
refusal---a run may contain several refusals and we do not count decisions---and we
report the model-driven and the rule-driven gate separately, because they refuse for
different reasons and at very different rates. We use ``disclosed'' only as shorthand for disclosure-marker positivity: the final answer matches the specified keyword list. This does not establish that the answer accurately explains the refusal.

\textbf{The instrument matters, so we report two.} The match is a keyword list, a weak
instrument: a refusal worded outside the list counts as undisclosed, so every count
here is a keyword-hit rate, not a validated semantic disclosure rate: the lists can miss an explicit disclosure or match unrelated inability wording. We report both lists rather than treating either as ground truth. List~A ($6$ markers) matches inability wording (\texttt{cannot},
\texttt{can not}, \texttt{unable}, \texttt{not able}, \texttt{outside the approved},
\texttt{denied}); list~B ($9$ markers) is an alternative list containing softer markers (\texttt{refus},
\texttt{confirm}, \texttt{alternative}, \texttt{instead}, \texttt{sorry}). For the
model-driven gate the disclosure count is $\mathbf{0}$ of $30$ refusal runs by
list~A and $\mathbf{9}$ of $30$ ($30.0\%$) by list~B under the 27B model, and $16$ of
$67$ ($23.9\%$) and $27$ of $67$ ($40.3\%$) under the 9B model. The gap between the
lists is itself informative: much of what counts as disclosure is detectable only
through a soft marker such as \texttt{sorry} or \texttt{instead}, not through an
explicit statement that the action was not performed.

\textbf{The rule-driven gate discloses more often, but not uniformly.}
\texttt{Baseline-detGate} discloses $12$ of $21$ refusal runs ($57.1\%$) by list~A and
$15$ of $21$ ($71.4\%$) by list~B under the 27B model, and $5$ of $16$ ($31.2\%$, both
lists) under the 9B model. That is higher than the model-driven gate under the 27B
model on both lists and under the 9B model on the stricter one, but \emph{lower} on
the 9B model's alternative list ($31.2\%$ against $40.3\%$), so we do not claim a uniform
advantage. These are descriptive marker comparisons between configurations with
different refusal sets; they do not identify why the final answers differ.

A possible failure mode worth studying is an omitted requested action whose omission
is not explained to the user. Our unvalidated keyword instruments do not establish
its semantic frequency. Section~\ref{sec:res-disclose-sens} instead tests how marker
counts change when the regeneration instruction is replaced and the refused-action
list is supplied.

\subsection{Cost, measured}
\label{sec:res-cost}

The reviewed version reported cost as an upper bound, because the audit loop breaks
at the first unambiguous verdict and the rounds actually consumed were not
recorded. Every model call now records its latency and, where the backend reports
them, its token counts, so Table~\ref{tab:cost-27b} (and
Table~\ref{tab:cost-qwen35_9b} for the second backbone) reports the result.

Three things follow, and each contradicts something the earlier version of this paper
said or implied.

\emph{The prompt-level policy is not free.} It spends no extra model calls---its
per-scenario call count is the planning call alone---but the safety sentence is paid on
every one of those calls. Against the ungated arm, which has neither a gate nor the
sentence, it adds $6278$ prompt tokens per $90$ cells under the 27B model
($+14.4\%$, about $70$ tokens per cell) and $6431$ ($+14.8\%$) under the 9B model.
``Zero extra cost'' is true of the call count and false of the token count.

\emph{The audit budget was a loose bound.} \texttt{Baseline-heavyAudit} triples the
round budget, from two rounds per dangerous action to three. Its realized cost is
indistinguishable from \texttt{Full}'s: $1.33$ against $1.33$ calls per scenario,
$0.27$ against $0.27$ rounds, and $11.6$s against $11.7$s of latency under the 27B
model; $1.38/1.39$ calls, $0.23/0.24$ rounds and $9.1/9.5$s under the 9B model. The
loop stops at the first verdict and verdicts arrive in the first round, so the third
round is never spent. Describing this arm as ``at least $3\times$ gate cost'', as the
earlier version did, overstated its cost by roughly that factor.

\emph{Call count does not order latency.} \texttt{Baseline-detGate} makes fewer calls
than \texttt{Full} ($1.23$ against $1.33$ per scenario under the 27B model) and no
auditor calls at all, yet its latency is \emph{higher}: $12.9$s against $11.7$s. The
same reversal holds under the 9B model ($1.18$ calls and $11.4$s against $1.39$ calls
and $9.5$s). Composition explains it: an audit reply is capped at $48$ tokens while a
regenerated answer is not, and a gate that refuses everything triggers regeneration far
more often. Forcing disclosure is the extreme case---\texttt{Fact-disclosePrompt} costs
$16.3$s per scenario against \texttt{Full}'s $11.7$s under the 27B model, a $39\%$
increase for one changed instruction, though only $4\%$ under the 9B model ($9.9$s
against $9.5$s).

We report latency as a \emph{cost} rather than a service-level metric: it is the sum
over every model call in the cell, so it grows with the number and length of calls
rather than describing a user-perceived response time.

\subsection{Memory eviction fires, but thinly}
\label{sec:res-memory}

The bounded memory (capacity 2) evicts $\mathbf{22}$ times across the nine
configurations under the 27B model and $\mathbf{23}$ times under the 9B model. The
forgetting ablation behaves as intended: \texttt{Abl-noForget} records $\mathbf{0}$
evictions under both backbones, which is a useful negative control on the mechanism
rather than a result about it.

The distribution is thin and lopsided. Under the 27B model six of the nine
configurations evict exactly twice, all in the same seed (per-seed $0/2/0$), while the
two configurations carrying the safety sentence evict five times ($2/0/3$). Under the
9B model the ungated arm evicts nine times ($0/6/3$) and the rest between one and
four. The reason is mechanical: memory is written only on successful scenarios and
success is $4$--$6\%$, so the pool fills slowly, and whether it overflows at all
depends on where the few successes happen to land in that seed's processing order.

We therefore still cannot report whether eviction helps or hurts. \texttt{Full} and
\texttt{Abl-noForget} are identical on every outcome we measure---success ($5/90$
both, under the 27B model), retained dangerous actions ($18$ both), auditor decisions
($24$ both, split $18/6$)---and they differ in exactly one respect: \texttt{Abl-noForget}
never evicts ($0$ against $2$). Two evictions across ninety runs is not a manipulation
of the mechanism; it is a reminder that the mechanism was barely switched on. This is a
negative result about the benchmark--model pair, not evidence that forgetting is
harmless.

\subsection{Exploratory cross-model replication}
\label{sec:res-crossmodel}

\textbf{All numbers in this subsection are exploratory, as are the nine-arm extensions on the primary backbone.} They come from a second
backbone run \emph{after} the primary 27B analysis, and they are reported
as a replication attempt rather than as confirmatory evidence; the deviation is
recorded in \texttt{evidence/analysis\_plan\_v4.md} (D1--D5).

We repeated the whole nine-configuration grid with a second recorded backbone
(Qwen3.5-9B; 27 group-seed files, 810 cells). The qualitative pattern of
\texttt{Full} versus \texttt{Baseline-promptSafety} replicates, the quantitative
picture moves, and part of it moves against the gate.

\paragraph{The prompt policy leaks under one backbone and not the other.}
\texttt{Full} leaves $0$ dangerous actions unreviewed under both models, and that zero
is now derived rather than assumed (\S\ref{sec:res-coverage}).
\texttt{Baseline-promptSafety} leaves $0$ under the 27B model but $\mathbf{2}$ under the
9B model, in a configuration that never constructs an auditor call. The two mechanisms
are therefore not interchangeable even where their counts coincide: one is a property of
the code path, the other a property of how reliably the model obeys an instruction, and
the second backbone is where that difference becomes visible.

\paragraph{On the second model the prompt arm scores higher.}
Under Qwen3.5-9B pooled success is $4/90$ ($4.4\%$) for \texttt{Full} and $\mathbf{8}/90$
($8.9\%$) for \texttt{Baseline-promptSafety}, against $5/90$ and $6/90$ under the 27B
model. As set out in \S\ref{sec:res-success} we do \emph{not} claim this as a gate cost:
at the scenario unit it is three informative units with exact $p=0.25$,
the cell-level $p=0.125$ does not survive Holm correction, and none of the four
discordant cells contains a gate decision of any kind. The one contrast we \emph{can}
attribute to the gate runs the other way and is described in that section: two
discordant cells in a single scenario, where the gate's denials cost the task.

\paragraph{The gate refuses far more often under the second model, and every refusal is
explicit.} The auditor's decisions differ sharply between backbones. Under the 27B model
\texttt{Full} approved $18$ and denied $6$ of the $24$ dangerous actions it saw ($25\%$
refusals); under the 9B model it approved $3$ and denied $19$ of $22$ ($86\%$). Each
decision now carries a verdict class, and across every gated configuration of both runs
there is \textbf{not one} decision classified \texttt{unparsed}, \texttt{all\_errors} or
\texttt{both\_tokens} (Tables~\ref{tab:audit-27b} and~\ref{tab:audit-qwen35_9b}). The
$86\%$ figure describes explicit auditor denials rather than parse-error defaults, which is
precisely what the earlier harness could not establish.

\paragraph{Marker contrasts across recorded backbones.} The alternative instruction
plus refused-action list increases list~B matches under both recorded backbones.
List~A rises only under 9B and remains zero under 27B. This is a marker contrast,
not a validated semantic-disclosure remedy (\S\ref{sec:res-disclose-sens}).

\paragraph{The forgetting mechanism's exercise is model-dependent.}
Under the 9B model eviction fires more often (\texttt{Abl-noNegot} $0/6/3$ per seed,
\texttt{Baseline-promptSafety} $0/2/2$) than under the 27B model ($0/2/0$ and $2/0/3$),
and \texttt{Abl-noForget} records none under either---the negative control behaving
correctly. Because memory is written only on success, how hard the forgetting mechanism
is exercised is governed by how often the model completes a task on this benchmark rather
than by the mechanism. Any claim that forgetting ``has no effect'' would therefore be a
claim about the model's success rate, not about the mechanism.

\subsection{Factorial ablation}
\label{sec:res-factorial}

Four of the configurations are the corners of a $2\times2$ over (auditor, safety
sentence) --- see Table~\ref{tab:design} --- so the two factors can be separated
rather than inferred from one contrast. Table~\ref{tab:factorial-27b} reports, at
the scenario unit, the cell means, the main effect of adding an auditor with the
prompt held fixed, the main effect of adding the safety sentence with the gate
held fixed, and their interaction, each with a paired percentile bootstrap
interval over scenarios.

Two qualifications bound what this design can isolate. Answer regeneration is not
an independent factor in the original arms: it fires only when the gate has refused
something, so it follows from having a gate rather than being a third factor to
cross. The \texttt{Fact-gateNoRegen} arm therefore isolates it by suppression
rather than by factorial crossing, and the difference between \texttt{Full} and
that arm attributes the regeneration step's effect on the answer, and hence on the
keyword-scored outcome. Separately, the deterministic arm is one fixed rule, so its
cell measures that rule rather than the class of policy-based enforcement it stands
in for.

Three readings follow from Table~\ref{tab:factorial-27b}, and all three hold in the
same direction under the second backbone.

\emph{Neither factor shows a detectable task-success effect.} The gate's main effect on success is
$+0.000$ with a bootstrap interval of $[0.000, 0.000]$ under the 27B model and
$-0.017$ $[-0.044, 0.000]$ under the 9B model; the safety sentence's is $+0.011$
$[-0.022, 0.044]$ and $+0.028$ $[-0.022, 0.083]$. Every interval contains zero, and so
does every interaction. These exploratory intervals do not establish a task-success effect or equivalence. Direct gate effects are difficult to resolve because few scenarios reach the gate: the model proposes a
dangerous action in $7$ of the $30$ under the 27B backbone (idx 2, 6, 7, 14, 22, 26, 29)
and in $8$ under the 9B backbone, so direct denial-related outcome differences can arise only in a small subset; prompt and memory effects could still change other scenarios.

\emph{Retained-action contrasts suggest diminishing marginal effects.} Under the 27B model the safety sentence's main effect is $-0.233$
$[-0.400, -0.083]$ while the gate's is $-0.033$ $[-0.083, 0.000]$: the pointwise interval for the marginal gate contrast includes zero, whereas that for the marginal safety-sentence contrast does not. These contrasts average across the other factor and do not establish comparative efficacy. The gate-alone raw contrast is $24$ to $18$ retained dangerous actions. Under
the 9B model the two are comparable, at $-0.128$ $[-0.239, -0.044]$ and $-0.117$
$[-0.233, -0.028]$. The interaction on this outcome is positive under both backbones
($+0.067$ and $+0.256$), consistent with diminishing marginal changes in this grid. These pointwise intervals are exploratory and do not establish a general substitution law.

\emph{Under the 27B model the substitution is complete.} On audited actions the gate's
main effect is $+0.133$ $[0.050, 0.233]$ and the sentence's is $-0.133$
$[-0.233, -0.050]$, with an interaction of $-0.267$ $[-0.467, -0.100]$: adding the
sentence to a gated system removes about as much auditing as adding the gate creates.
The cell values make it concrete. With the gate alone the 27B model audits $0.267$
dangerous actions per scenario; with the gate and the sentence together it audits
$\mathbf{0.000}$, because the model no longer proposes any. In that cell the gate is
not merely redundant---it is unreachable.

This is the cleanest statement the design supports about the two mechanisms: they act
at different points in the pipeline. The safety sentence acts on \emph{proposal} and the
gate acts on \emph{retention}, so a sufficiently compliant planner leaves the gate
nothing to do. It also bounds what our auditability contrast can claim, because a gate
that never fires produces an empty audit log for the same reason a prompt policy does,
and only the configuration table tells the two apart.

\subsection{How sensitive is disclosure to the instruction?}
\label{sec:res-disclose-sens}

The gate's refusals often lack the tested disclosure markers to the user
(\S\ref{sec:res-disclosure}). The reviewer asked whether that is a property of the
mechanism or of the instruction we happened to give the regeneration step: if the
prompt explicitly requires refused actions to be named, does disclosure rise?
\texttt{Fact-disclosePrompt} tests exactly this, holding the gate and the
regeneration step fixed and changing only the regeneration instruction, which is
additionally handed the refused list.

It rises, and by a wide margin---but the size of the rise depends on which keyword list
scores it and on which backbone produces it. Under the 9B model \texttt{Full} discloses
$2$ of $13$ refusal runs ($15.4\%$) on both lists, while \texttt{Fact-disclosePrompt}
discloses $\mathbf{8}$ of $14$ ($57.1\%$) on the strict list and $\mathbf{14}$ of $14$
($100\%$) on the alternative list. Under the 27B model \texttt{Full} discloses $0$ of $6$ on
both lists, and the disclosure arm $0$ of $6$ on the strict list but $\mathbf{6}$ of $6$
($100\%$) on the alternative list.

The two lists produce different descriptive counts. The strict list matches inability wording (\texttt{cannot},
\texttt{unable}, \texttt{denied}); the alternative nine-marker list includes \texttt{refus},
\texttt{sorry}, \texttt{instead}, \texttt{alternative} and \texttt{confirm}. Our
disclosure instruction asks the model to say which actions ``were refused'', and the
model complies using exactly that stem---which the strict list does not contain. Under
the 27B model the regenerated answers apparently avoid the strict list's phrasings
altogether, so that instrument scores no improvement whatever even as the alternative list
scores every run.

We therefore answer the question conditionally. Replacing the regeneration instruction and supplying the refused-action list changes marker counts from $0/6$ to $6/6$ under the 27B model and from $2/13$ to $14/14$
under the 9B model, on the list~B. But whether a given answer counts as
disclosed is partly a property of the instrument, and a stricter instrument can show no
change at all. The honest summary is that disclosure is both instruction-sensitive and
instrument-sensitive, and neither sensitivity is visible unless both arms and both lists
are reported.

The instruction's cost is measurable and appears in \S\ref{sec:res-cost}: $39\%$ more
latency per scenario under the 27B model and $4\%$ under the 9B model, with task success
unchanged ($5/90$ and $4/90$, identical to \texttt{Full}).

A second arm points the same way from the other side. \texttt{Fact-gateNoRegen} keeps
the gate and its denials but skips regeneration, leaving the original plan answer in
place---and it discloses \emph{more} than \texttt{Full} on the list~B: $3$ of $6$
against $0$ of $6$ under the 27B model, and $7$ of $14$ against $2$ of $13$ under the 9B
model. Our original regeneration prompt excludes actions outside the approved list.
The alternative arm changes this instruction and additionally supplies the refused-action
list. These comparisons support sensitivity of marker counts to the tested regeneration
choices, rather than identifying the effect of a single sentence or measuring semantic
explanation quality. On list~A, removing regeneration changes nothing under the 27B
model ($0$ of $6$ either way).

We also report how sensitive the disclosure count is to the measuring instrument.
Two keyword lists are used across our analysis scripts (a six-token and a
nine-token variant), and a refusal worded outside both is counted as undisclosed.
The count is therefore a marker-based measurement rather than a semantic bound on disclosure, and the two lists do \emph{not} agree: on the model-driven gate they give $0$ of $30$
refusal runs by the strict list versus $9$ of $30$ by the alternative list at 27B, and
$16$ of $67$ versus $27$ of $67$ at 9B (\texttt{\_lists\_agree}: \texttt{false} in both
analysis records), so each count is quoted together with the list that produced it.
A disclosure that neither
list catches would still be scored as silent, which is why we present the
instruction's effect as a contrast between arms rather than as an absolute rate.

\section{Discussion}
\label{sec:discussion}

\paragraph{A coverage metric cannot see the mechanism.}
The result we consider most useful is a negative one about measurement. The standard
quantity one would report for an approval gate---how many dangerous actions were retained
without confirmation---reads zero for every gated configuration in our harness, and it read
zero for the system-prompt policy under the primary backbone too. In an earlier version of
this work that zero was \emph{assigned} rather than measured, which is why it could not
fail; it is now derived from the retained plan and the audit log, and the same expression
returns $24$ and $25$ for the ungated arms and $2$ for the prompt policy under the weaker
backbone. Deriving it makes the metric honest but does not make it discriminating: a gate
that reviews every dangerous action, and a prompt policy the model happens to obey, both
read zero, and the difference between them is invisible in the number a practitioner would
report. Once we count dangerous actions \emph{retained} in the final plan instead, the two
mechanisms separate ($18$ versus $0$ under one backbone, $3$ versus $2$ under the other),
and they separate further on the audit record ($24$ recorded decisions versus none). A
practitioner choosing between them is choosing between auditability and cost, and the usual
scalar is blind to exactly that difference.

\paragraph{Token cost and review coverage are different quantities.}
The prompt policy cost zero extra model calls but $14.4\%$ more prompt tokens, and under
the primary backbone it never left a dangerous action unconfirmed; under the weaker
backbone it leaked two. Its coverage is a property of the model's compliance rather than of
the configuration, so the present records do not establish deployment reliability for either mechanism; the derived zero verifies consistency of retained plans with logs, not enforcement at a real execution boundary. The factorial exposes
a further problem: under the primary backbone, adding the safety sentence to a gated system
leaves the gate with nothing to audit, so the two are substitutes rather than complements,
and a deployment should measure whether its auditor is reached before attributing outcomes to it. We also
cannot detect over-safety here: the benchmark carries no safety labels, so a
refusal that saves the user money and a refusal that blocks a legitimate request
are scored identically.

\paragraph{Refusal markers depend on the regeneration design.}
Across both backbones the model-driven gated configurations contain $97$ refusal runs;
$16$ final answers match list~A and $36$ match list~B. These runs include correlated
configurations and scenario repeats, so the aggregate is descriptive. Replacing the
regeneration instruction and supplying the refused-action list changes list~B counts
from $0/6$ to $6/6$ under 27B and from $2/13$ to $14/14$ under 9B. The latter
comparison also has a different refusal denominator. Skipping regeneration gives
another marker contrast. These records motivate semantic evaluation of refusal
explanations; they do not establish the frequency or cause of unexplained refusals in
deployed systems. Observed latency changes of $39\%$ and $4\%$ are specific to the
recorded serving stacks and run conditions.

\paragraph{The forgetting mechanism was barely exercised.}
Our capacity-2 memory evicted $22$ times across the nine configurations of the 27B grid
and $23$ times under the 9B model, and six of the nine 27B configurations evicted only
twice, all within the same seed. Because memory is written only on success and success is
$4$--$6\%$, the pool rarely fills. We therefore report a failure of experimental
stimulation rather than a null effect: on a benchmark where the model rarely completes a
task, no conclusion about eviction policies can be drawn. That eviction fires more often
under the second model suggests its exercise is coupled to per-model competence on this
benchmark rather than to the mechanism itself, and \texttt{Abl-noForget} recording zero
evictions under both backbones confirms the ablation does what it claims.

\paragraph{Three measurements need different evidence.}
Coverage asks whether every designated retained action has a recorded review. Decision correctness asks whether approval or denial agrees with an explicit authorization policy. Benefit asks whether the intervention improves outcomes under a suitable counterfactual comparison. Our records support descriptive coverage accounting; they lack authorization ground truth and provide weak task-success evidence. These three questions should not be answered by a single zero-valued counter.

\section{Limitations}
\label{sec:limitations}

\paragraph{Thin danger basis.} Only a minority of scenarios contain a dangerous
action; the gate-vs-no-gate contrast is carried by 7 of the 30 scenario units,
and the refusals concentrate in one action family. The gate's behaviour on the
other whitelisted actions is untested.

\paragraph{No ground-truth safety labels.} The benchmark's assertions decide only
whether the task completed; they do not say whether an action was safe. We
therefore cannot score a refusal as correct or as over-rejection, and every
statement about safety in this paper is a statement about \emph{coverage}, not
about harm avoided. Zhang et al.~\cite{zhang2026guard} argue this boundary is a
matter of authorization policy, which is why we do not claim the gate makes the
system safer.

\paragraph{Self-audit.} The auditor is the same base model as the executor, with
a short prompt. We do not compare it with an independent auditor and do not measure
false-accept or false-reject rates for it.

\paragraph{Unchecked assertions.} $12.1\%$ of our cell-assertions ($432$ of $3564$) are
never evaluated, because they name no known tool action. The figure is identical across
all nine configurations and both backbones, but missing assertions can still interact with arm-specific behavior; equal missingness does not establish an unbiased success comparison. Making the prefix
dispatch case-insensitive recovered three assertions per configuration that an earlier
version of the scorer silently skipped; the remainder are semantic assertions that no
keyword scorer of this kind can check.

\paragraph{What the factorial settles, and what it does not.} An earlier version of this
work compared a gate against a prompt policy while varying three things at once and could
not attribute the difference. The $2\times2$ over (auditor, safety sentence) now
separates those two factors, and \texttt{Fact-gateNoRegen} isolates answer regeneration.
Two things remain unisolated. Regeneration is not a third factorial factor, because it
fires only when the gate has refused something, so it sits downstream of the gate rather
than orthogonal to it; we isolate it by suppression, not by crossing. And
\texttt{Baseline-detGate} is one fixed rule, so its cell measures that rule rather than
the class of policy-based enforcement it stands in for---a real implementation would
carry a policy language, argument-level predicates and a proof obligation, none of which
we built.

\paragraph{Weak scorer for user assertions.} Assertions that do not name a tool action
are \emph{not scored at all}: they are recorded as unchecked (see \S\ref{sec:setup}).
Assertions carrying a \texttt{user:} prefix are scored by a keyword-overlap heuristic, a
weak instrument whose token list is given in \S\ref{sec:res-disclosure}. This limits the
absolute success numbers, which are low ($5.6\%$ under the 27B model and $4.4\%$ under
the 9B model for the gated configurations), and a low success rate in turn removes the
power needed to detect any effect on task success.

\paragraph{Cost is measured, but not decomposed in the manuscript.} The harness records
the rounds each audit consumed and the latency and token counts of every call, so the
strengthened gate's cost is a measurement rather than a bound---and the measurement shows
the bound was loose, since tripling the round budget changed nothing observable. What the
manuscript does not break out is the per-phase split (planning against auditing against
regeneration); those counts are in the delivered JSON but not in the tables. Latency is
also wall-clock on a shared GPU, so it carries contention noise we did not control.

\paragraph{Models and domain.} Both models are from the same family (Qwen) and we
use a single domain (travel). The cross-model agreement we report is therefore
agreement within a family, not across families.

\paragraph{Non-determinism.} Decoding is not deterministic even at temperature
zero, so the three seeds change processing order rather than constituting
independent draws; we treat the 30 scenarios as the independent unit and could
not fully remove pseudo-replication at the cell level.

\section{Conclusion}
\label{sec:conclusion}

We compared an approval gate with a system-prompt policy and a deterministic rule
in a plan-level travel harness: $30$ scenarios, three processing orders, nine
configurations, two recorded model-service labels, and $1620$ cells. Retained per-cell
records support comparisons of proposed and retained actions, recorded review,
marker-positive answers, and serving-stack costs.

The count of retained dangerous actions without recorded approval is derived from
plans and logs. It is zero in all gated configurations, $24$ and $25$ in the ungated
arms, and $0$ and $2$ in the prompt-policy arms. This consistency measure does not
establish execution-boundary enforcement or correct authorization. Counting retained
actions separately reveals a different contrast: Full versus prompt retains $18$
versus $0$ under 27B and $3$ versus $2$ under 9B. Full records $24$ and $22$ review
decisions; the deterministic rule also records review, using no auditor-model calls.
Its total latency is not uniformly lower.

In the 27B factorial grid, adding the safety sentence removes the dangerous proposals
that would otherwise reach the gate. The exploratory interaction on audited actions
is $-0.267$, with a pointwise $95\%$ interval of $[-0.467,-0.100]$. No task-success
effect is established at the scenario unit; this dataset does not resolve competence
trade-offs or demonstrate equivalence. The 9B Full--prompt contrast ($4/90$ versus
$8/90$) rests on three informative scenario units, with exact sign-test $p=0.25$.

Marker-positive answers depend on the tested regeneration choices and keyword list.
The alternative instruction plus refused-action list increases list~B matches from
$0/6$ to $6/6$ under 27B and from $2/13$ to $14/14$ under 9B, with observed
stack-specific latency increases. Semantic explanation quality and deployment effects
remain unmeasured.

Finally, the harness has no tool-execution layer, no authorization ground truth, and
low benchmark success. Parse failures affect $16.7\%$ of 9B Full cells, compared with
$5.6\%$ in its prompt arm, and can make an empty plan appear compliant. The defensible
result is therefore an exploratory measurement case study: audit trails, retained
plans, parsing, explanation proxies, and cost must be reported separately. We do not
establish a task-outcome or harm-prevention benefit in this dataset.

\section*{Broader Impact}

This paper measures a common control pattern in a plan-level harness, and
its main practical message is cautionary in two directions. First, a coverage
count---the number of unconfirmed dangerous actions---should not be used on its
own to justify a runtime control, because a system-prompt instruction can
match it on one backbone while leaving no record of what was considered or refused; teams that
need an audit trail should measure the audit trail. Second, the control we
studied can omit a requested action from the retained plan. Only 16 of 97 model-gated refusal runs matched list~A and 36 matched list~B; these unvalidated proxies motivate checking whether users receive an accurate explanation. They do not establish that the remaining answers were silent or harmful. We do not claim that approval gates make a
system safer, in either direction: the benchmark we use has no ground-truth
labels for whether an action was safe, so a refusal is indistinguishable from an
over-rejection in our measurements, and a policy that refuses more is scored the
same as one that refuses correctly. We also note that the auditor in our
configuration is the same model as the executor, which is a weak self-audit
setting; deployments should not take our numbers as evidence that
same-model review is adequate. No human subjects, personal data, or
high-stakes decisions are involved; the benchmark scenarios are authored
enterprise examples and the models are run locally on a single shared GPU.

\bibliographystyle{tmlr}
\bibliography{main}

\end{document}
